
\subsubsection{2.1 Execution Feedback for Code LLMs}

CodeRL introduced using program execution outcomes as reward signals for code generation. RLTF extended this by combining multiple feedback granularities: coarse (pass/fail), fine-grained (token-level penalties at error locations), and adaptive (based on test case outcomes). Their ablation studies showed that combined approaches outperform single-signal methods.

RLEF incorporated textual execution feedback directly into model prompts during training, modifying the input rather than the reward structure.

This work builds on RLTF's multi-granularity framework but identifies that uniform application of fine-grained feedback ignores localization reliability.

\subsubsection{2.2 Aggregation and Credit Assignment}

VeRPO identified cardinality bias---where rewards computed over different numbers of test cases create artificial variance---and proposed variance-reduced aggregation. Their analysis demonstrates that naive dense reward application can degrade rather than help training.

This work addresses a complementary dimension: localization reliability rather than aggregation. VeRPO asks ''how should we combine signals?''; we ask ''which signals should we apply where?''

\subsubsection{2.3 Error Localization in Programming}

The program repair and fault localization literature has recognized that tracebacks vary in usefulness. Spectrum-based fault localization methods weight suspicious lines based on test outcome correlations rather than trusting stack traces directly. However, this insight has not been incorporated into RL reward design for code generation.

RLTF's error categorization (U\_line, U\_global, U\_ignore) implicitly acknowledges this variation but does not exploit it for conditional feedback application.

\subsubsection{2.4 Credit Assignment in RL}

The credit assignment problem---determining which actions led to which outcomes---is fundamental to RL. This work addresses credit assignment across tokens. When traceback localization is unreliable, fine-grained feedback misassigns credit to wrong tokens.

